\section{Methods}

\subsection{Reference data and force-free backbone}

We use the flexible SPC/Fw water model \citep{spcfw} with 64 molecules (192 atoms) in a periodic
cubic cell. The NVE training data, the NVT reference, and the independent NVT trajectory are
described in the main text. The reference NVT trajectory is 2.5 ns at 300 K with a CSVR
thermostat \citep{csvr}, a 50 fs thermostat damping time, and a 0.5 fs saving interval; the
readout is fitted on a 100 ps block starting at 500 ps (20,000 frames at a 5 fs stride). The fit
uses only the mode velocity ACF. The RDF, VDOS, and diffusion are not used to fit any readout
parameter and are held out as evaluation observables. As a block-sensitivity check, we refitted
the AR(10) coefficients on the 0--100 ps block: the coefficients agree with the 500--600 ps fit
to \(r=0.999\) (maximum coefficient difference 0.08), with prediction-error ratios
0.0044/0.058/0.107 versus 0.0042/0.058/0.115 for translation/rotation/internal motion. The fit is
therefore not sensitive to the choice of reference block. The
independent NVT trajectory uses a different initial condition and seed: 100 ps of thermalization
followed by 500 ps of production at 300 K, with the first production frame at 301.8 K. The
reference data are read-only; no reference force is used during propagation.

The backbone is a relational equivariant message-passing network with a 6 \AA{} cutoff,
192 scalar channels, 48 vector channels, 96 edge channels, 16 radial channels, and a 5 fs
horizon. It is trained on 5 fs pairs from the NVE trajectories (5,000 training, 1,250 validation,
and 1,000 test graphs) by minimizing a normalized endpoint position and velocity loss, with batch
size 8, a learning rate of \(3\times10^{-4}\), warmup--cosine scheduling, gradient clipping at 10,
and 250,000 optimizer steps. The selected checkpoint is epoch 398 (2,357,616 parameters). The
published TrajCast water checkpoint \citep{trajcast} has 2,258,856 parameters, uses the native e3nn
backend \citep{e3nn} and a 5 fs horizon, and is evaluated under the same 100 ps protocol with its
native CSVR thermostat. In the GLR experiments the backbone is frozen: no backbone parameter is updated, and
the readout is fitted separately.

\subsection{Mode decomposition}

For each molecule \(m\), let \(M_m=\sum_{i\in m}m_i\), \(R_m=M_m^{-1}\sum_{i\in m}m_i r_i\), and
\(V_m=M_m^{-1}\sum_{i\in m}m_i v_i\). With \(r_{mi}=r_i-R_m\) (minimum image) and
\(v_{mi}=v_i-V_m\), the angular velocity is
\begin{equation}
\omega_m=I_m^{-1}\sum_{i\in m}m_i\,r_{mi}\times v_{mi},\qquad
I_m=\sum_{i\in m}m_i\left(|r_{mi}|^2\mathbf 1-r_{mi}r_{mi}^{\mathsf T}\right),
\end{equation}
using a pseudo-inverse for linear molecules. The three mass-weighted modes are
\begin{equation}
v_i^{\rm trans}=V_m,\qquad
v_i^{\rm rot}=\omega_m\times r_{mi},\qquad
v_i^{\rm int}=v_{mi}-v_i^{\rm rot}.
\end{equation}
They are orthogonal under the mass-weighted inner product, so
\(K=\sum_a K_a\) with \(K_a=\tfrac12\sum_i m_i|v_i^a|^2\). The mode degrees of freedom are
\(N_{\rm trans}=3N_{\rm mol}-3\), \(N_{\rm rot}=3N_{\rm mol}\), and
\(N_{\rm int}=3N_{\rm atom}-6N_{\rm mol}\).

\subsection{Mode-resolved AR($p$) readout}

For each mode \(a\), the reference velocity autocorrelation is
\begin{equation}
\alpha_a(\ell)=\frac{\sum_i m_i\langle v_i^a(t)\cdot v_i^a(t+\ell\Delta t)\rangle}
{\sum_i m_i\langle |v_i^a(t)|^2\rangle}.
\end{equation}
We fit an AR(\(p\)) process by Yule--Walker,
\begin{equation}
v^a(t)=\sum_{k=1}^{p}c_k^a v^a(t-k\Delta t)+\varepsilon^a(t),\qquad
\mathbf R_a\mathbf c_a=\mathbf r_a,
\end{equation}
where \((\mathbf R_a)_{ij}=\alpha_a(|i-j|)\) and \((\mathbf r_a)_k=\alpha_a(k)\). The
prediction-error ratio is
\begin{equation}
\rho_a=1-\sum_{k=1}^{p}c_k^a\alpha_a(k)
=\operatorname{Var}(\varepsilon^a)/\operatorname{Var}(v^a).
\end{equation}
The fluctuation--dissipation condition fixes the noise variance,
\begin{equation}
\operatorname{Var}(\varepsilon^a)=\rho_a\,N_a k_BT/2.
\end{equation}
At each step a fresh Gaussian field is decomposed into the three modes and each mode component is
rescaled to this target noise kinetic energy before it is added to the AR prediction. The main
configuration uses \(p=10\); the order ablation uses \(p=1,5,20\).

\subsection{Readout, blend, and variance rescaling}

The readout velocity is
\begin{equation}
\tilde v_t=\sum_a\left(\sum_{k=1}^{p}c_k^a v_{t-k}^a+\varepsilon_t^a\right).
\end{equation}
It is blended with the backbone velocity,
\begin{equation}
v_t^{\rm blend}=(1-\lambda)v_t^{\rm pred}+\lambda\tilde v_t,
\end{equation}
with \(\lambda=0.1\), and each mode is rescaled to its canonical mean kinetic energy:
\begin{equation}
v_t^{\rm GLR}=\sum_a s_a\,(v_t^{\rm blend})^a,\qquad
s_a=\sqrt{\frac{N_a k_BT/2}
{\tfrac12\sum_i m_i|(v_t^{\rm blend})_i^a|^2}}.
\end{equation}
Equation~(11) uses the canonical mean. In the stochastic configuration, the group kinetic energy
is drawn from \(K_a\sim\Gamma(N_a/2,k_BT)\) and the group is rescaled to that target; the
deterministic rescale is the mean-only limit and is reported as an ablation. The stochastic target
restores the canonical kinetic-energy variance. The water 2.5 ns production uses the deterministic
mean target because the stochastic variant restores the KE variance (0.437/0.429 versus 0.440 eV)
but gives a Green--Kubo integral (\(1.16\times10^{-9}\)) inconsistent with the MSD
(\(1.44\times10^{-9}\)). Quartz uses the stochastic target (0.398 versus 0.4017 eV) because the
deterministic species rescale suppresses the total KE variance to 0.0036 eV. The rescale
partition \(a\) is not
necessarily the memory-mode basis. Water uses translation, rotation, and internal modes; quartz
uses five phonon bands (0--5, 5--10, 10--20, 20--30, and 30--40 THz), with a coarser Si/O species
partition as a control. The partition must contain enough degrees of freedom that the Gamma scale
factor is close to one; a per-mode rescale with one degree of freedom per group gives an
\(O(1)\) random scale every step and destroys the mode phase coherence (Supplementary
Information). The quartz configuration is the phonon-band stochastic rescale with the memory
kernel switched off; the phonon memory kernel is a negative control because the 5 fs backbone's
mode ACF is already correct.

The total centre-of-mass velocity is projected out after the rescale. The position update remains
\(q_{t+1}=q_t+\Delta q_t\), where \(\Delta q_t\) is the backbone displacement; the readout is not
integrated as a second position update. We tested the alternative
\(q_{t+1}=q_t+\Delta q_t+\Delta t(v_t^{\rm GLR}-v_t^{\rm pred})\) and partial
translation-only and molecule-COM variants in the ablation. The full correction is unstable, and
the partial variants do not close the diffusion gap without breaking the MSD--Green--Kubo
consistency or the RDF (see main text and Supplementary Information).

\subsection{Rollout protocol and metrics}

All rollouts use a 5 fs horizon, no external thermostat, and the same safety stop as the
production screens (minimum interatomic distance 0.65 \AA, pre-thermostat temperature 3000 K).
The 100 ps screens (20,000 steps) use the val\_04 NVE frame and the independent NVT frame as
initial conditions; the 2.5 ns production runs (500,000 steps) start from the same frames with
no thermalization. The seeds are 2026110401 (val\_04), 2026110402 (independent), and
2026110401/02/03 for the seed study. The published checkpoint is run with its native CSVR
thermostat under the same 100 ps protocol; its 2.5 ns diffusion is taken from the published value
because the native float64 rollout is prohibitively slow on our contended GPU.

The mode kinetic partition is the mean \(K_a\) over the production window. The ACF is computed
from the mode velocities with the same mass-weighted estimator as the fit, at lags 5--1000 fs.
RDFs are computed on 512 uniformly spaced frames with a 6 \AA{} maximum radius. VDOS overlaps use
the same block spectral estimator as the reference: six blocks of 4.096 ps, a 100 THz Nyquist
frequency, and a Hann window. The diffusion coefficient is obtained from the molecular
centre-of-mass mean-squared displacement with time-origin averaging, eight blocks, a 50 fs
sampling cadence, and a linear fit over 10--100 ps per block; the global centre-of-mass drift is
removed. For 2.5 ns, each block is 312.5 ps long. The diffusion coefficient is
\(D=\mathrm{slope}/6\times10^{-8}\) m\(^2\) s\(^{-1}\) when the MSD is in \AA\(^2\) and the lag in
ps. Error bars are block means and standard deviations; the seed study uses three independent
noise realizations.

The GLR rollout costs 15.4 ns day\(^{-1}\) end-to-end on an NVIDIA RTX PRO 6000 Blackwell GPU for
100 ps, with a peak memory of 185 MB and 2.36M parameters. The readout adds a mode decomposition
and a \(p\)-step history per step; it is negligible compared with the backbone forward pass.

\subsection{Baselines and controls}

The structure-only control is the 250k backbone without the readout. FDR v1 is the AR(1) readout
with the same blend and rescale. M0 is the 250k backbone fine-tuned on 5 fs pairs from the
reference NVT trajectory. The CSRV control is the 5-step moment-loss rollout. The native ep145
checkpoint is a partially trained TrajCast architecture and is reported as a diagnostic, not a
fair architecture baseline. The published TrajCast checkpoint is the external baseline.

\subsection{Statistics and reproducibility}

The reference partition has substantial block-to-block scatter: the first 100 ps block gives
2.373/2.430/2.609 eV, while the 2.5 ns mean is 2.459/2.502/2.450 eV. We therefore use the 2.5 ns
mean for long-time partition comparisons and report the 100 ps block separately. All error bars
are block means and standard deviations over eight contiguous blocks unless stated otherwise. The
diffusion fits use the same eight-block protocol and report the block mean and standard
deviation; the fit quality is summarized by \(R^2\). The seed study repeats the full 100 ps
rollout with independent noise realizations and reports the seed-to-seed spread of the partition,
RDF, VDOS, and diffusion. The reference trajectories are read-only; the rollout code writes a new
HDF5 file and refuses to overwrite an existing one. Each rollout records the checkpoint SHA-256,
the seed, the initial frame, the integration step, and the safety-stop status. No reference force
is used during propagation; the reference forces are used only in the independent energy
verification of the NVT trajectory.

\subsection{Computational cost}

The GLR rollout costs 15.4 ns day\(^{-1}\) end-to-end on an NVIDIA RTX PRO 6000 Blackwell GPU for
a 100 ps run, with a peak memory of 185 MB. The readout adds a mode decomposition and a \(p\)-step
history per step; the cost is negligible compared with the backbone forward pass. The published
native checkpoint is about an order of magnitude slower in our environment because it uses
float64 and the native e3nn path, which is why its 2.5 ns baseline is taken from the published
value rather than re-run.

\subsection{Data and code availability}

The reference NVT trajectory, the independent NVT trajectory, and the rollout trajectories are
available at [repository/DOI to be supplied]. The published TrajCast checkpoint is available from
the IBM Hugging Face repository \texttt{ibm-research/trajcast.models-arxiv2025}. The GLR readout,
the Yule--Walker fit, and the rollout code are available at [repository URL to be supplied]. No
reference force is used during propagation.
